\chapter{Exponentials and Logarithms}
\label{ch:explog}

\begin{tcolorbox}[feynbox={Sky}{BBg},title={Before and After}]
\textbf{Before this chapter you need:} percentages and gross returns
(\S\ref{sec:percentages}), functions and their domains
(\S\ref{sec:functions}--\S\ref{sec:domain}), composition
(\S\ref{sec:composition}).\\
\textbf{After it you will understand:} why every quantity in this thesis is a
\emph{log} return rather than a percentage return, why that choice makes
multi-day returns add up exactly, and why a log-scaled histogram is the honest
way to look at the tail of a distribution.
\end{tcolorbox}

This chapter earns its length. The decision to work in log returns propagates
into every model, every metric and every figure in the project, and the reasons
are worth seeing in full rather than accepting on authority.

%% ─────────────────────────────────────────────────────────────────────────────
\section{Compound Growth}
\label{sec:compound}

\situation{
You put \$100 in an account paying 8\% a year, and you leave the interest in.
After one year you have \$108. The second year's 8\% is charged on \$108, not on
the original \$100, so you gain \$8.64 rather than \$8. Growth feeds on itself.
}

\problem{
Because each year's gain depends on the running total, growth over $n$ years is
not $n$ separate additions. It is $n$ multiplications, and multiplication is far
harder to reason about than addition --- especially when you want to compare
periods of different length, or split a long period into pieces.
}

\workedarith{
\$100 at 8\% per year, interest retained.

\medskip
\begin{tabular}{@{}llll@{}}
\toprule
Year & Start & Gain & End \\
\midrule
$1$ & $100.00$ & $8.00$ & $108.00$ \\
$2$ & $108.00$ & $8.64$ & $116.64$ \\
$3$ & $116.64$ & $9.33$ & $125.97$ \\
\bottomrule
\end{tabular}

\medskip
\textbf{As repeated multiplication:}
$100 \times 1.08 \times 1.08 \times 1.08 = 100 \times 1.08^3 = 125.97$. \checkmark

\medskip
\textbf{Compare to simple (non-compounding) interest:} $100 + 3 \times 8 = 124$.
The gap of \$1.97 after only three years is the compounding effect, and it widens
with time.

\medskip
\textbf{Now compound more often.} Take 100\% per year, but split it into $n$
instalments of $1/n$ each:
\[
\begin{aligned}
  n = 1:&\quad (1 + 1/1)^1 = 2.00000 \\
  n = 2:&\quad (1 + 1/2)^2 = 1.5^2 = 2.25000 \\
  n = 4:&\quad (1 + 1/4)^4 = 1.25^4 = 2.44141 \\
  n = 12:&\quad (1 + 1/12)^{12} = 2.61304 \\
  n = 365:&\quad (1 + 1/365)^{365} = 2.71457
\end{aligned}
\]
The numbers climb, but ever more slowly, and they are converging on something
just above $2.718$. They do not run away to infinity.
}

\formaldef{
An \textbf{exponential function} has the form $f(x) = a^x$ with $a > 0$. Growth
at rate $R$ per period, compounded over $n$ periods, multiplies the starting
value by $(1 + R)^n$.

The limit reached by the table above defines \textbf{Euler's number}:
\[
  e = \lim_{n \to \infty}\left(1 + \frac{1}{n}\right)^n \approx 2.718281828\ldots
\]
It is irrational: its decimal expansion never terminates or repeats. The function
$f(x) = e^x$ is called \emph{the} exponential function, written $\exp(x)$.
}

\analogybreak{
The bank analogy suggests $e$ arises from a peculiar accounting convention, as if
it were an artefact of how often interest is credited. It is not. The quantity
$e^x$ is singled out by a property with nothing to do with money: it is the only
function that is its own rate of change, a fact Chapter~\ref{ch:calculus1} makes
precise. Continuous compounding is one place $e$ shows up, not the reason it
exists.
}

\whyhere{
Continuous compounding is the natural setting for asset prices, which are quoted
continuously through the trading day rather than at fixed interest dates. Writing
$P_1 = P_0 e^{r}$ instead of $P_1 = P_0(1 + R)$ is what makes $r$ additive across
time, which is the property the rest of this thesis is built on.
}

%% ─────────────────────────────────────────────────────────────────────────────
\section{The Logarithm as the Inverse}
\label{sec:loginverse}

\situation{
A machine doubles whatever you feed it: put in 3, get 6; put in 5, get 10. Now
someone hands you a 32 and asks what was fed in. You need the machine run
backwards. Since $2^5 = 32$, the answer is 5 --- you are asking ``how many
doublings?''.
}

\problem{
Running exponentials backwards is exactly the question ``what power was this?''.
It cannot be answered by any combination of the four arithmetic operations, so it
needs a name and a function of its own. Without one, every compound-growth
calculation is a one-way street.
}

\workedarith{
Base 2, so we are counting doublings.

\medskip
\begin{tabular}{@{}llll@{}}
\toprule
$x$ & $2^x$ & & $\log_2$ of it \\
\midrule
$0$ & $1$  & & $\log_2 1 = 0$ \\
$1$ & $2$  & & $\log_2 2 = 1$ \\
$2$ & $4$  & & $\log_2 4 = 2$ \\
$3$ & $8$  & & $\log_2 8 = 3$ \\
$5$ & $32$ & & $\log_2 32 = 5$ \\
\bottomrule
\end{tabular}

\medskip
Read the table left to right for exponentials, right to left for logarithms. The
two columns hold identical information.

\medskip
\textbf{The characteristic move.} Note that $4 \times 8 = 32$, and in the log
column $2 + 3 = 5$. Multiplying the numbers corresponds to adding their
logarithms. This is not a coincidence of these particular values; it is the
defining property, and it is the entire reason logarithms are useful.

\medskip
\textbf{With base $e$.} The natural logarithm $\ln$ undoes $e^x$:
$\ln(1) = 0$, $\ln(e) = 1$, $\ln(e^2) = 2$. Numerically
$\ln(2) \approx 0.693$ and $\ln(10) \approx 2.303$.
}

\formaldef{
The \textbf{logarithm to base $a$}, written $\log_a$, is the inverse of $x \mapsto a^x$:
\[
  \log_a(y) = x \iff a^x = y .
\]
The \textbf{natural logarithm} $\ln = \log_e$ satisfies
\[
  e^{\ln x} = x \quad\text{and}\quad \ln(e^x) = x .
\]
Its domain is the strictly positive reals: $\ln(0)$ and $\ln(-1)$ are undefined,
because no real power of $e$ produces zero or a negative number. As
$x \to 0^{+}$, $\ln x \to -\infty$.
}

\analogybreak{
Calling $\ln$ ``the inverse'' invites the assumption that it undoes $e^x$ in every
sense, including numerically. It does not. Computing $e^{-40}$ gives roughly
$4 \times 10^{-18}$, which in ordinary floating-point arithmetic may round to
something whose logarithm is no longer $-40$. The mathematical inverse is exact;
the computational one is not, which is why numerical code so often carries
quantities in log form from the start rather than converting back and forth.
}

\whyhere{
The strictly-positive domain is the reason we take logarithms of \emph{prices}
and never of returns. A price is always positive; a return is negative on roughly
half of all days, and $\ln$ has nothing to say about those.
}

%% ─────────────────────────────────────────────────────────────────────────────
\section{The Log Rules}
\label{sec:logrules}

\situation{
Before electronic calculators, engineers multiplied large numbers by looking up
two logarithms in a printed table, adding them by hand, and looking the result
back up. A hard operation was traded for an easy one at the cost of two lookups.
}

\problem{
The same trade is worth making today, for a different reason. Products of many
terms overflow or underflow the range a computer can represent, and they resist
algebraic manipulation. Sums do neither. Anything that converts products to sums
buys both numerical stability and analytical tractability.
}

\workedarith{
Verify each rule on numbers small enough to check mentally, using base 2.

\medskip
\textbf{Product rule.} $\log_2(4 \times 8) = \log_2 32 = 5$, and
$\log_2 4 + \log_2 8 = 2 + 3 = 5$. \checkmark

\medskip
\textbf{Quotient rule.} $\log_2(32 / 4) = \log_2 8 = 3$, and
$\log_2 32 - \log_2 4 = 5 - 2 = 3$. \checkmark

\medskip
\textbf{Power rule.} $\log_2(4^3) = \log_2 64 = 6$, and
$3 \times \log_2 4 = 3 \times 2 = 6$. \checkmark

\medskip
\textbf{Log of 1.} $\log_2 1 = 0$, because $2^0 = 1$. This holds in every base,
and it is the fact that makes ``no change'' correspond to a log return of exactly
zero.
}

\formaldef{
For $x, y > 0$ and any real $k$:
\[
  \ln(xy) = \ln x + \ln y, \qquad
  \ln\!\left(\frac{x}{y}\right) = \ln x - \ln y, \qquad
  \ln(x^k) = k \ln x, \qquad
  \ln 1 = 0 .
\]
To change base: $\log_a x = \ln x / \ln a$.
}

\analogybreak{
There is no rule for $\ln(x + y)$. The logarithm converts multiplication into
addition, not addition into anything simpler --- and beginners routinely invent
$\ln(x+y) = \ln x + \ln y$, which is false. Check it: $\ln(2 + 2) = \ln 4 \approx 1.386$,
while $\ln 2 + \ln 2 \approx 1.386$. Those agree, which is exactly the sort of
accident that entrenches the error. Try $\ln(1 + 3) = \ln 4 \approx 1.386$ against
$\ln 1 + \ln 3 = 0 + 1.099 = 1.099$. They differ.
}

\whyhere{
The product rule is the engine of the next section. It is also why the
log-likelihood, rather than the likelihood, is what gets maximised whenever a
GARCH model is fitted to a return series in this project: a product of thousands
of small probabilities underflows to zero, while the corresponding sum of
logarithms is perfectly well behaved.
}

%% ─────────────────────────────────────────────────────────────────────────────
\section{Why Log Returns Are Additive}
\label{sec:logreturns}

\situation{
You invest \$100. It rises 10\%, then falls 10\%. You are back where you started
--- or so the symmetry of the two numbers suggests. You are not. You have \$99.
}

\problem{
Percentage returns do not add across time. Over two days the errors are small
enough to ignore; over 5{,}000 trading days --- the span of this project's data
--- they compound into nonsense. Any quantity that is to be summed, averaged, or
modelled as a sequence of increments must be additive by construction.
}

\workedarith{
\textbf{The asymmetry, in full.} Start at 100.
\[
  100 \times 1.10 = 110, \qquad 110 \times 0.90 = 99 .
\]
Net change $-1\%$, not $0\%$. In gross terms $1.10 \times 0.90 = 0.99$.

\medskip
Now in logs:
\[
  \ln(1.10) = +0.09531, \qquad \ln(0.90) = -0.10536 .
\]
Sum: $+0.09531 - 0.10536 = -0.01005$. And directly,
$\ln(99/100) = \ln(0.99) = -0.01005$. \checkmark

The log returns sum to exactly the right answer. The percentages do not.

\medskip
\textbf{The genuinely symmetric case.} Take a rise and fall that really do cancel:
up by a factor $1.10$, then down by a factor $1/1.10 = 0.909091$.
\[
  \ln(1.10) = +0.09531, \qquad \ln(0.909091) = -0.09531 .
\]
They sum to exactly zero, and indeed $100 \times 1.10 \times 0.909091 = 100$.
In log space, ``up then back down'' is $+a$ then $-a$ --- perfectly symmetric.
In percentage space it is $+10\%$ then $-9.09\%$, which looks lopsided but is not.

\medskip
\textbf{Three days from the data-shaped example.}
Prices 4{,}000, 4{,}080, 4{,}050.
\[
  r_1 = \ln(4080/4000) = \ln(1.02) = +0.019803
\]
\[
  r_2 = \ln(4050/4080) = \ln(0.992647) = -0.007380
\]
Sum: $+0.012423$. Direct: $\ln(4050/4000) = \ln(1.0125) = +0.012423$. \checkmark

By percentages: $+2.000\% - 0.735\% = +1.265\%$, whereas the true change is
$+1.250\%$. An error of $1.5$ basis points over two days --- and it accumulates.
}

\formaldef{
The \textbf{log return} on day $t$ is
\[
  r_t = \ln\!\left(\frac{P_t}{P_{t-1}}\right) = \ln P_t - \ln P_{t-1} .
\]
Summing over $T$ days telescopes:
\[
  \sum_{t=1}^{T} r_t = (\ln P_1 - \ln P_0) + (\ln P_2 - \ln P_1) + \cdots + (\ln P_T - \ln P_{T-1})
  = \ln P_T - \ln P_0 = \ln\!\left(\frac{P_T}{P_0}\right) ,
\]
every interior term cancelling with its neighbour. The relationship to the simple
return $R_t$ is $r_t = \ln(1 + R_t)$, and for small moves $r_t \approx R_t$: at
$R = 0.02$ we get $r = 0.019803$, agreeing to three decimal places.
}

\analogybreak{
Log returns are additive across \emph{time} but not across \emph{assets}. The log
return of a portfolio is not the weighted sum of its constituents' log returns,
because a sum of prices does not decompose under a logarithm --- there is no rule
for $\ln(x+y)$, as \S\ref{sec:logrules} insisted. Portfolio work therefore often
returns to simple returns, which \emph{are} additive across assets and not across
time. Each convention buys one form of additivity and forfeits the other.
}

\whyhere{
Every model, metric and figure in this thesis consumes $r_t$. Prices are never
used directly: they trend, they live on incomparable scales across the five BRICS
markets, and they violate the stationarity that Chapter~\ref{ch:timeseries}
requires. Converting to log returns fixes all three problems in one step.
}

%% ─────────────────────────────────────────────────────────────────────────────
\section{Reading Tails on a Log Scale}
\label{sec:logscale}

\situation{
A histogram of daily returns is a spike at zero with two barely visible smears at
the edges. The bar at zero is thousands of days tall; the bar at $-8\%$ holds
three days. Drawn on ordinary axes, the three-day bar is invisible --- less than
one pixel --- and the picture says only ``returns are usually small'', which
everybody knew.
}

\problem{
The rare days are the ones that matter. A risk model is judged almost entirely on
how it handles the far edges of the distribution, and those edges are exactly what
a linear axis erases. We need a way of drawing the picture in which a count of 3
and a count of 3{,}000 are both legible.
}

\workedarith{
Suppose four bins hold counts $3{,}000$, $300$, $30$, $3$.

\medskip
\textbf{On a linear axis}, with the tallest bar 10\,cm high, the others are
$1$\,cm, $1$\,mm and $0.1$\,mm. The last is thinner than the ink.

\medskip
\textbf{On a log axis}, we plot $\log_{10}$ of each count:
\[
  \log_{10} 3000 = 3.477, \quad
  \log_{10} 300 = 2.477, \quad
  \log_{10} 30 = 1.477, \quad
  \log_{10} 3 = 0.477 .
\]
The heights are now $3.477, 2.477, 1.477, 0.477$ --- equally spaced, one unit
apart. Every bar is visible, and the constant gap of exactly $1$ reveals that each
bin holds one tenth of the previous, a tenfold decay that the linear plot
concealed completely.
}

\formaldef{
A \textbf{logarithmic axis} places a value $v$ at position $\log v$. Equal
distances then represent equal \emph{ratios} rather than equal differences: the
gap from 1 to 10 is the same length as from 10 to 100.

A distribution whose tail decays as a power law, $\Prob(|X| > x) \sim x^{-\alpha}$,
becomes a straight line of slope $-\alpha$ when both axes are logarithmic, since
\[
  \ln \Prob(|X| > x) = -\alpha \ln x + \text{constant} .
\]
Straightness on a log--log plot is the visual signature of a power law, and the
slope reads off the tail index directly.
}

\analogybreak{
A log axis cannot display a count of zero, since $\ln 0$ is undefined --- empty
bins simply vanish rather than sitting on the axis. It also flatters the eye:
differences that are genuinely enormous are compressed into a modest vertical
gap, so a model that is off by a factor of ten can look close. The log axis makes
the tail \emph{visible}; it does not make disagreements there \emph{small}.
}

\whyhere{
The stylised-facts figure produced for every market in this project includes a
log-scale distribution panel for exactly this reason, with values outside the
plotted range clipped rather than dropped so that no observation silently
disappears. The panel is where a generator that has matched the centre of the
distribution but missed its tails is caught --- and Chapter~\ref{ch:heavytails}
shows that missing the tails is the characteristic failure of these models.
}
